\section{Foundation Geometry Model Benchmark}
\label{sec:benchmark}

To select the best backbone for our calibration pipeline, we benchmark seven recent foundation geometry models on our dataset. All models receive the same 8 synchronised camera images per timestep and must output camera poses without any prior information.

\subsection{Experimental Protocol}

We sample 91 timesteps at regular intervals from synchronised video (2688$\times$1520 resolution, resized per model requirements). For each model, we:
\begin{enumerate}
    \item Run inference on all timesteps (8 cameras per timestep)
    \item Aggregate per-timestep poses via quaternion averaging (rotations) and median (translations)
    \item Apply Sim(3) Procrustes alignment to ground truth (with reflection check)
    \item Compute ATE, ARE, RPE, and co-located pair distances
\end{enumerate}

Ground truth poses are obtained from COLMAP applied to a static scene capture (no cattle present). Two camera pairs (C2--C3 and C5--C6) share physical mounting poles, providing a structural validation metric (ground truth distance $\approx$ 0).

\subsection{Results}

\begin{table}[t]
\centering
\caption{Foundation geometry model benchmark on our wide-baseline, feature-adversarial dataset. Models are ranked by ATE. ``Type'' indicates whether the model processes all views jointly (N-view) or in pairs with post-hoc global alignment. All metrics use Sim(3)-aligned predictions; co-location distances (C2C3, C5C6) measure structural consistency. Pairwise models (MonST3R, Pow3R) use only 20 timesteps due to computational cost.}
\label{tab:benchmark}
\begin{tabular}{@{}llccccccc@{}}
\toprule
Model & Venue & Type & ATE$\downarrow$ & ARE$\downarrow$ & RPE$_R$$\downarrow$ & RPE$_t$$\downarrow$ & C2C3$\downarrow$ & Time (s) \\
\midrule
\textbf{Pi3X} & CVPR'25 & N-view & \textbf{0.177} & \textbf{4.37\degree} & \textbf{6.26\degree} & \textbf{0.39} & \textbf{0.088} & 106 \\
MapAnything & CVPR'25 & N-view & 0.627 & 9.17\degree & 14.19\degree & 0.93 & 0.313 & 86 \\
VGGT & CVPR'25 & N-view & 1.248 & 10.48\degree & 12.82\degree & 1.55 & 2.011 & 75 \\
Pow3R & CVPR'25 & Pairwise & 1.835 & 180\degree & 28.53\degree & 7.24 & 0.218 & 497 \\
NoPo4D & arXiv'25 & N-view & 2.556 & 62.06\degree & 77.01\degree & 3.17 & 0.627 & 214 \\
MonST3R & ICLR'25 & Pairwise & 2.607 & 47.24\degree & 23.87\degree & 3.56 & 0.111 & 1492 \\
Fast3R & CVPR'25 & N-view & 2.743 & 96.64\degree & 102.88\degree & 4.36 & 3.183 & 460 \\
\bottomrule
\end{tabular}
\end{table}

\subsection{Analysis}

The benchmark reveals a clear hierarchy:

\paragraph{Tier 1: Pi3X.} Achieves sub-5\degree{} rotation error and 0.177 ATE---3.5$\times$ better than the next model. Its permutation-equivariant architecture that processes all 8 views jointly is well-suited to wide-baseline multi-camera geometry.

\paragraph{Tier 2: MapAnything, VGGT.} Both achieve $\sim$10\degree{} rotation error. Usable but significantly less accurate than Pi3X on this task.

\paragraph{Tier 3: Remaining models.} MonST3R, Pow3R, NoPo4D, and Fast3R all exceed 20\degree{} rotation error---effectively failing on this dataset. The pairwise models (MonST3R, Pow3R) struggle because their post-hoc global alignment amplifies errors when inter-camera overlap is minimal.

\paragraph{Key observations.}
\begin{itemize}
    \item N-view joint models (Pi3X, MapAnything, VGGT) consistently outperform pairwise models, confirming that joint multi-view reasoning is essential for wide-baseline configurations.
    \item The co-location metric (C2C3) discriminates structural understanding: Pi3X correctly identifies that C2 and C3 share a pole (0.088), while VGGT (2.011) and Fast3R (3.183) place them far apart.
    \item All models were published in 2025, indicating that wide-baseline feature-adversarial calibration remains a frontier challenge for even the latest foundation models.
\end{itemize}

Based on this benchmark, we select Pi3X as our backbone and develop our self-refinement method (Sections~\ref{sec:stage2}--\ref{sec:stage4}) to push beyond its $\sim$4.4\degree{} accuracy ceiling.
